% =============================================================================
% bab3-metode.tex — BAB III METODE PENELITIAN
% Skripsi / Tesis / Disertasi (Lampiran I)
% =============================================================================

\chapter{METODE PENELITIAN}
\label{bab:metode}

\section{Jenis Penelitian}
\label{sec:jenis-penelitian}

Penelitian ini merupakan penelitian [kuantitatif/kualitatif/mixed methods]. Jelaskan
jenis penelitian dan pendekatan yang digunakan beserta alasannya.

\section{Variabel Operasional}
\label{sec:variabel-operasional}

Pada penelitian ini, variabel-variabel yang digunakan adalah sebagai berikut:

\begin{table}[H]
\caption{Variabel Operasional}
\label{tab:variabel-operasional}
\begin{isitable}
\begin{tabular}{p{0.5cm} p{2.5cm} p{3.5cm} p{3.5cm} p{2.5cm} p{1.5cm}}
\toprule
\textbf{No} & \textbf{Variabel} & \textbf{Definisi Operasional} & \textbf{Indikator} & \textbf{Skala} & \textbf{Sumber} \\
\midrule
1 & Variabel X & Definisi operasional variabel X & Indikator 1, 2, 3 & Likert & Sumber \\
2 & Variabel Y & Definisi operasional variabel Y & Indikator 1, 2 & Likert & Sumber \\
\bottomrule
\end{tabular}
\end{isitable}
\end{table}

\section{Tahapan Penelitian}
\label{sec:tahapan-penelitian}

Penelitian ini dilaksanakan melalui tahapan-tahapan berikut:

\begin{enumerate}
  \item Identifikasi masalah dan penetapan topik penelitian.
  \item Studi literatur dan pengembangan kerangka teoritis.
  \item Penyusunan instrumen penelitian.
  \item Pengumpulan data.
  \item Pengolahan dan analisis data.
  \item Penarikan kesimpulan dan penyusunan laporan.
\end{enumerate}

\section{Populasi dan Sampel}
\label{sec:populasi-sampel}

\subsection{Populasi}
\label{subsec:populasi}

Populasi dalam penelitian ini adalah [uraikan populasi penelitian]. Jumlah populasi
dalam penelitian ini adalah sebanyak \ldots

\subsection{Sampel}
\label{subsec:sampel}

Pengambilan sampel dilakukan dengan menggunakan teknik [purposive/random/stratified]
sampling. Penentuan jumlah sampel menggunakan rumus Slovin sebagai berikut:

\begin{equation}
n = \frac{N}{1 + N \cdot e^2}
\label{eq:slovin}
\end{equation}

\noindent
di mana $n$ = jumlah sampel, $N$ = jumlah populasi, dan $e$ = batas toleransi
kesalahan (5\%).

\section{Pengumpulan Data dan Sumber Data}
\label{sec:pengumpulan-data}

\subsection{Sumber Data}
Data yang digunakan dalam penelitian ini terdiri dari:
\begin{enumerate}
  \item Data primer: data yang diperoleh langsung dari responden melalui kuesioner.
  \item Data sekunder: data yang diperoleh dari sumber-sumber seperti laporan tahunan,
        publikasi ilmiah, dan sumber lainnya.
\end{enumerate}

\subsection{Teknik Pengumpulan Data}
Teknik pengumpulan data yang digunakan adalah:
\begin{enumerate}
  \item Kuesioner (angket): diberikan kepada responden secara langsung/daring.
  \item Studi dokumentasi: analisis dokumen-dokumen yang relevan.
\end{enumerate}

\section{Uji Validitas dan Reliabilitas}
\label{sec:uji-validitas-reliabilitas}

\subsection{Uji Validitas}
Uji validitas dilakukan untuk mengukur ketepatan instrumen penelitian. Instrumen
dikatakan valid apabila nilai korelasi $r_{hitung} > r_{tabel}$ pada taraf signifikansi
5\%.

\subsection{Uji Reliabilitas}
Uji reliabilitas dilakukan dengan menggunakan koefisien Cronbach's Alpha ($\alpha$).
Instrumen dinyatakan reliabel apabila nilai $\alpha \geq 0.60$.

\section{Teknik Analisis Data}
\label{sec:teknik-analisis}

\subsection{Analisis Deskriptif}
Analisis deskriptif digunakan untuk menggambarkan karakteristik responden dan
distribusi jawaban responden terhadap setiap indikator.

\subsection{Analisis [Regresi/SEM/dsb.]}
Teknik analisis yang digunakan dalam penelitian ini adalah [nama teknik]. Uraikan
alasan pemilihan teknik analisis dan model persamaannya:

\begin{equation}
Y = \alpha + \beta_1 X_1 + \beta_2 X_2 + \varepsilon
\label{eq:model-regresi}
\end{equation}

\noindent
di mana $Y$ = variabel dependen, $X_1$ dan $X_2$ = variabel independen, $\alpha$ =
konstanta, $\beta_1$ dan $\beta_2$ = koefisien regresi, $\varepsilon$ = error term.
